\clearpage\begingroup\SourceLeaf{372b}\ArithmeticHead{41}{R}{DIE GRUNDLAGEN DER GEOMETRIE.}
\begin{center}
\textbf{9.}\par\smallskip
\textls[100]{Nachwirkung der \textsc{Gauss}schen Gedanken.}
\end{center}
\noindent\hspace*{\parindent}Bei der Zurückhaltung, die sich \textsc{Gauss} zur Regel gemacht hatte, haben
während seines Lebens nur wenige Bevorzugte etwas von seinen Ansichten
über die Grundlagen der Geometrie erfahren, und die Eingeweihten haben
ihr Wissen für sich behalten. Zum Beispiel hat \textsc{Dirichlet}, mit dem \textsc{Gauss}
bei dessen Besuch im März 1827 von der nichteuklidischen Geometrie gesprochen
hatte (W. VIII, S. 188), untersucht, wie sich die Potentialtheorie im
nichteuklidischen Raume gestalte, aber nichts darüber veröffentlicht (\textsc{Lob.} S. 444).
In weiteren Kreisen wurde erst etwas davon bekannt, als \textsc{Sartorius} 1856 in
seiner Schrift \textit{Gauss zum Gedächtniss} berichtete, \textsc{Gauss} habe eine selbständige
Geometrie ausgebildet, die gelte, wenn man das Parallelenaxiom nicht zugebe
(W. VIII, S. 267—268). Diese Andeutung wurde bald darauf bestätigt durch
den 1860 herausgekommenen zweiten Band des Briefwechsels zwischen \textsc{Gauss}
und \textsc{Schumacher} (Briefe vom Jahre 1831, W. VIII, S. 210—219), und 1865
erschien der fünfte Band mit dem Briefe vom 28. November 1846 (W. VIII,
S. 238), durch den die Aufmerksamkeit auf \textsc{Lobatschefskij} gelenkt wurde.
Nachdem jetzt, um mit \textsc{Hoüel} zu reden, die imposante Autorität \textsc{Gaussens}
\par\vfill\noindent{\footnotesize X\textsubscript{2} Abh. 4.\hfill 6\hspace*{2em}}
\clearpage\SourceLeaf{373}\ArithmeticHead{42}{L}{STÄCKEL, GAUSS ALS GEOMETER.}
\noindent gesprochen hatte, fand der Hinweis auf J. \textsc{Bolyai} und \textsc{Lobatschefskij} Beachtung,
den \textsc{Baltzer} 1867 in der zweiten Auflage seiner \textit{Elemente der
Mathematik} gab; durch ihn angeregt veröffentlichte \textsc{Hoüel} französische Übersetzungen
der \textit{Geometrischen Untersuchungen} und der \textit{Scientia spatii absolute vera}
und machte so diese verschollenen Schriften allgemein zugänglich. Damit
war der Boden vorbereitet für eine verständnisvolle Aufnahme der zu derselben
Zeit aus \textsc{Riemanns} Nachlass herausgegebenen Habilitationsrede vom
Jahre 1854: \textit{Über die Hypothesen, welche der Geometrie zu Grunde liegen};
dazu kamen 1868 die Aufsätze von \textsc{Helmholtz}. Während man bis dahin die
Beschäftigung mit dem elften Axiom als ein Vorrecht unklarer Köpfe angesehen
und mit den Bemühungen um die Quadratur des Kreises und das Perpetuum
mobile auf eine Stufe gestellt hatte, erregten jetzt die Untersuchungen
über die Grundlagen der Geometrie allgemeine Teilnahme, und als später
noch die Kritik der Arithmetik hinzukam, entstand ein neuer Zweig der
Mathematik, der als Axiomatik bezeichnet wird.
\clearpage\endgroup
